% ---------------------------------------------------------------------------
% NasAge validation figures and captions
%
% Two separate stories, deliberately kept apart:
%   Figure 5 / S5  -- mucociliary ageing: does the manuscript's own ageing signal
%                     replicate in independent nasal cohorts?
%   Figure 6 / S6  -- the mucosal programme panel and its relation to infection.
%
% Compile with:  pdflatex figures.tex
% The PDFs live beside this file; no graphics path is needed.
% ---------------------------------------------------------------------------
\documentclass[11pt,a4paper]{article}

\usepackage[utf8]{inputenc}
\usepackage[T1]{fontenc}
\usepackage[margin=25mm]{geometry}
\usepackage{graphicx}
\usepackage{caption}
\usepackage{amsmath}
\usepackage{microtype}
\usepackage[hidelinks]{hyperref}

\captionsetup{font=small,labelfont=bf,justification=justified,singlelinecheck=false,
              width=\textwidth,skip=8pt}
\setlength{\parindent}{0pt}
\setlength{\parskip}{6pt}

\renewcommand{\figurename}{Figure}

\title{NasAge validation: Figures 5 and 6}
\author{}
\date{Generated 25 November 2026}

\begin{document}
\maketitle

\noindent
\textbf{Scope.} Every panel below is drawn from a results file produced by an analysis plan
committed to version control \emph{before} the corresponding data were accessed. No panel
recomputes a model statistic; the plotting scripts read results and arrange them. Deviations
from each locked plan are dated in the corresponding \texttt{PREREG\_RESULTS.md}.

\textbf{The two figures are separate stories and share no panel.} Figure 5 asks whether the
manuscript's mucociliary ageing signal replicates in independent nasal cohorts. Figure 6 asks
whether a panel of curated mucosal programmes is cell-type specific, changes with age, and
relates to infection. The programme-specificity analysis appears in Figure 6 only.

\clearpage

% ===========================================================================
% FIGURE 5
% ===========================================================================
\setcounter{figure}{4}
\begin{figure}[t]
  \centering
  \includegraphics[width=\textwidth]{figure5.pdf}
  \caption{\textbf{The mucociliary ageing signal replicates in an independent nasal cohort.}
  \textbf{(a)} M02 mucociliary axoneme score against donor age in the replication cohort
  ($n=120$ donors, bulk nasal brushing). Line and band are the ordinary least-squares fit
  and its 95\% confidence interval. The Spearman $\rho$ is computed from the plotted points;
  the adjusted $\beta$ is the age coefficient from the pre-registered model, which includes
  sequencing depth as a covariate. Benjamini--Hochberg correction is over the declared family
  of 13 modules, in which this test does not survive ($\mathrm{BH}=0.22$).
  \textbf{(b)} The same analysis for CILIUM28, the gene set derived in the primary cohort.
  Because CILIUM28 was defined on the primary data, this panel is a transfer test, not an
  independent derivation.
  \textbf{(c)} Both age coefficients against the distribution of 500 expression- and
  size-matched random gene sets scored identically in the same cohort. Neither observed
  coefficient is approached by any random set, and \textbf{0 of 500} random sets reach
  $p<0.05$ against a nominal 5\%, so the association is specific to these gene sets rather
  than a property of any set of comparably expressed genes.
  \textbf{(d)} Interferon-stimulated gene score (M05) in children versus adults in Loske
  2022, by epithelial cell type, as donor means. Horizontal bars are group medians.
  Benjamini--Hochberg values are within the declared family; \emph{emp.\ p} is the
  expression-matched empirical $p$, and an asterisk marks the combinations that clear it.
  The ciliated contrast reaches $\mathrm{BH}=1\times10^{-3}$ but fails its matched null
  ($\mathrm{emp.}\ p=0.076$) and is therefore not claimed.
  \textbf{(e)} The same contrast for M06 MHC-I processing, which clears both the family-wise
  correction and the matched null in all three cell types.
  \textbf{Interpretation.} The childhood-to-adult contrast is the robust part of this
  signal. Within adulthood the gradient is flat (Supplementary Figure S5a), so panels d and e
  should not be read as a continuous decline across the lifespan.}
  \label{fig:cilium}
\end{figure}

\clearpage

% ===========================================================================
% FIGURE 6
% ===========================================================================
\begin{figure}[t]
  \centering
  \includegraphics[width=\textwidth]{figure6.pdf}
  \caption{\textbf{A cell-type-specific mucosal programme panel, and its relation to
  infection.}
  \textbf{(a)} Mean score of each curated programme in each epithelial family, across three
  single-cell cohorts and 224 donor $\times$ cell-type units. The solid box marks the cell
  type each programme was curated for; the dashed orange box marks the one programme (P6
  viral entry factors) that scores highest somewhere else.
  \textbf{(b)} Standardised mean difference between a programme's own cell type and the
  other two, with 95\% confidence intervals. Five of six programmes are specific and the sign
  is unanimous in every cohort. P6 is significantly \emph{lower} in its assigned cell type
  ($\mathrm{SMD}=-1.66$) and is reported as misassigned rather than dropped.
  Note that curated cell-type genes are expected to be enriched in that cell type, so the
  magnitude of these SMDs overstates the achievement; the informative result is the one
  programme that fails.
  \textbf{(c)} M06 antigen presentation against age, 4--77 years, as donor $\times$ cell-type
  means with per-family least-squares fits. The decline is present in all three cell types
  ($\mathrm{BH}\le1\times10^{-4}$), is robust across 12 analytic specifications, and clears
  its expression-matched null under all four scoring methods in basal and secretory cells
  (three of four in ciliated). The slope is carried by the childhood-to-adult contrast:
  within adulthood it is flat in all three cohorts tested.
  \textbf{(d)} Odds ratio per standard deviation of baseline P5 mucin barrier score for
  infection outcome, in three challenge cohorts, with a fixed-effect inverse-variance pooled
  estimate. The direction is protective in all three. Only pneumococcal carriage clears its
  own expression-matched null (panel f); RSV and NOVA do not, so the pooled estimate is shown
  and not claimed.
  \textbf{(e)} Specification curve for P5 in the RSV and pneumococcal cohorts: 16
  specifications each, being four scoring methods $\times$ four covariate sets, ordered by
  estimate. Filled points are $p<0.05$. All 32 specifications fall below $\mathrm{OR}=1$
  (100\% sign stability in both cohorts). Sign stability of this kind is evidence that a
  multiplicity correction can neither supply nor remove.
  \textbf{(f)} The observed P5 effect in pneumococcal carriage against 500 expression- and
  size-matched random gene sets scored identically. The empirical $p$ is 0.000 under every
  one of the four scoring methods. By a rule fixed before any null was drawn, the reported
  figure for each cohort is the \emph{worst} scorer, never the best.
  \textbf{Limits that govern the whole figure.} No vaccine outcome exists in any reachable
  cohort, so nothing here supports a claim about vaccine responsiveness; a live attenuated
  influenza vaccine does not engage any of these programmes despite a working positive
  control (Supplementary Figure S6d). Cohort sizes are 39, 39 and 74. P5's outcome
  association does not survive if all 140 primary tests in the project are treated as one
  family.}
  \label{fig:mvri}
\end{figure}

\clearpage

% ===========================================================================
% SUPPLEMENTARY S5
% ===========================================================================
\renewcommand{\thefigure}{S5}
\begin{figure}[t]
  \centering
  \includegraphics[width=\textwidth]{figure5_supp.pdf}
  \caption{\textbf{Boundary conditions on the mucociliary ageing result.}
  \textbf{(a)} The interferon score shows no gradient with age \emph{within} older adults, in
  two independent cohorts ($\rho=-0.07$, $p=0.46$, $n=120$; $\rho=-0.16$, $p=0.32$, $n=39$).
  Together with Figure 5d this localises the signal to the childhood-to-adult transition and
  excludes a continuous adult decline.
  \textbf{(b)} Detected genes per cell in CDKN2A$^+$ versus CDKN2A$^-$ cells. p16$^+$ status
  tracks sequencing depth almost entirely: the medians differ by 1{,}377 genes per cell and
  181 of 187 donors show the same direction. Senescence calls from droplet data are therefore
  confounded with capture efficiency, which is why the corresponding cohort is withdrawn from
  biological claims.
  \textbf{(c)} Fraction of each module's genes detected in at least 5\% of cells, in four
  cohorts and two chemistries. The dashed line is the pre-specified 60\% measurability gate.
  M12 SASP falls at or near zero everywhere and is unmeasurable in droplet data; it is
  retired from primary analysis rather than reported as a null.
  \textbf{(d)} The three outcome-anchored tests with their pre-registered directions,
  verdicts and the reason for each verdict. Predictions, directions and multiplicity were
  fixed by a committed plan before the data were accessed.}
  \label{fig:s5}
\end{figure}

\clearpage

% ===========================================================================
% SUPPLEMENTARY S6
% ===========================================================================
\renewcommand{\thefigure}{S6}
\begin{figure}[t]
  \centering
  \includegraphics[width=\textwidth]{figure6_supp.pdf}
  \caption{\textbf{Everything the programme panel failed.}
  This figure exists so that Figure 6 is not read as a complete account.
  \textbf{(a)} Evidence matrix: for each of five candidate axes, whether it is measurable,
  cell-type specific, changes with age, relates to infection, and is engaged by a vaccine.
  Filled green is a pass, small gold is partial, orange is a failure, a dash is not tested.
  P5's infection entry is \emph{partial}, not a pass, because it clears its matched null in
  one cohort of three.
  \textbf{(b)} Declared axes versus the effective number of independent tests (Li \& Ji,
  2005) in each cell type. Five declared axes behave as three to four, because PC1 carries
  55--69\% of their variance, so Benjamini--Hochberg over the declared count is the wrong
  correction in both directions.
  \textbf{(c)} Age slopes for all five axes in all three cell types. A green marker denotes
  the combinations that clear the expression-matched null. Seven of eleven tested
  combinations reach $\mathrm{BH}<0.05$ yet fail that null: a general age drift that most
  gene sets partially capture. Only antigen presentation is robust.
  \textbf{(d)} Paired change in each programme after live attenuated influenza vaccine. The
  M05 interferon positive control rises as expected, confirming the assay works on these
  samples, while every mucosal programme is null. Three independent designs returned this
  result. \textbf{This is the panel that prevents the figure set being titled as a validated
  vaccine-readiness index.}
  \textbf{(e)} Median pairwise correlation among the five modules that panel v2 collapsed
  into one interferon axis, in the two cohorts the collapse was derived from and in three
  held-out cohorts. Black rules mark the 95th percentile of random five-module subsets of the
  same panel. The axis is 0.72--0.93 where it was derived and 0.34--0.50 where it was not,
  beating chance in one held-out cohort of three. By the pre-registered rule, the collapse
  does not replicate and v1 remains the panel of record.
  \textbf{(f)} Empirical $p$ for every programme against its expression- and size-matched
  null, in each cohort, under each of four scoring methods (mean $z$, median $z$, rank,
  PC1 of set). Boxed cells clear the null. Specificity is scorer-dependent for most
  programmes, which is why the reported figure is the worst scorer rather than the best.
  \textbf{(g)} Fraction of \emph{random} expression-matched gene sets that reach model
  $p<0.05$, against the nominal 0.05. The bulk challenge cohorts are close to nominal
  (0.02--0.06), but the single-cell pseudobulk path is not (0.33 basal, 0.43 secretory,
  0.76 ciliated). Where this fraction is high, a nominal $p$ from any gene set is
  anticonservative and the matched null is the only valid test. This is the dominant-axis
  problem measured directly, and it is the reason the trajectory results in panel c are
  judged on their matched nulls rather than on their $p$-values.}
  \label{fig:s6}
\end{figure}

\end{document}
